\section*{第 \ref{ch:b3:stable-isotopes} 章——稳定同位素地球化学}

\begin{solution}{exo:b3:stable-isotopes:1}
$\alpha = 1030.5/997.9 = 1.032669$；$1000\ln\alpha = 32.15$；$\Delta = 30.5 -
(-2.1) = \qty{32.6}{\permil}$。二者相差 \qty{0.45}{\permil}，
即命题中的 $(\delta_A^2 - \delta_B^2)/2000 = 0.46$。
\end{solution}

\begin{solution}{exo:b3:stable-isotopes:2}
$\delta^{18}\mathrm O = 0.69\times(-25) - 13.6 = \qty{-30.9}{\permil}$；
$\delta\mathrm D = 5.6\times(-25) - 100 = \qty{-240}{\permil}$；$d = -240 +
8\times30.85 = \qty{6.8}{\permil}$。
\end{solution}

\begin{solution}{exo:b3:stable-isotopes:3}
$\Delta_B = -26.0 - 1.0 = \qty{-27.0}{\permil}$；$f_{\mathrm{org}} = (-5 -
1.0)/(-27.0) = 0.22$。
\end{solution}

\begin{solution}{exo:b3:stable-isotopes:4}
$T^2 = 0.0391\times10^6/(0.600 - 0.154) = \num{87668}$，$T = \qty{296.1}{K}$，
即 \qty{22.9}{\celsius}。
\end{solution}

\begin{solution}{exo:b3:stable-isotopes:5}
$f^{0.01} = 967/987$，$f = \exp(100\ln0.97974) = 0.129$：\qty{87}{\percent} 的
水汽已经降落。雨水：$1.0100\times967 - 1000 =
\qty{-23.3}{\permil}$。
\end{solution}

\begin{solution}{exo:b3:stable-isotopes:6}
$\alpha = 1000/960 = 1.041667$；$1000\ln\alpha = 40.82$；$1000(\alpha - 1) =
41.67$；$\Delta = 40$。两种近似的误差都约为 \qty{0.8}{\permil}，
其中第二种的误差是第一种的两倍：在 \qty{40}{\permil} 处，
一阶展开已经不够精确。
\end{solution}

\begin{solution}{exo:b3:stable-isotopes:7}
$\delta_c = (1.50 + 29.99)/0.97001 = \qty{32.46}{\permil}$（VSMOW）；$1000\ln\alpha
= 1000\ln1.032464 = 31.95$；$T = 18030/(31.95 + 32.42) = \qty{280.1}{K}$，
即 \qty{7.0}{\celsius}。分馏每开尔文约变化 \qty{0.23}{\permil}，
因此水的 \qty{0.1}{\permil} 变化约相当于 \qty{0.4}{\celsius}。
\end{solution}

\begin{solution}{exo:b3:stable-isotopes:8}
冰量贡献 $125\times0.011 = \qty{1.375}{\permil}$；其余部分为 $1.77 - 1.375 =
\qty{0.40}{\permil}$。在 \qty{2}{\celsius} 附近，斜率为 $18030/275.15^2 =
\qty{0.238}{\permil.K^{-1}}$，因此深层海水
约冷了 \qty{1.7}{\celsius}。
\end{solution}

\begin{solution}{exo:b3:stable-isotopes:9}
$\Delta^{33}\mathrm S = 8.2 - 0.515\times10.0 = \qty{3.05}{\permil}$，远离 0：
非质量相关硫同位素，只有在没有臭氧屏蔽时才会形成。黄铁矿
形成于大气中氧含量不显著的时期，时间早于约
\qty{2.45}{\Ga}。
\end{solution}

\begin{solution}{exo:b3:stable-isotopes:10}
在封闭系统中，重同位素守恒：$R_0 = fR_v + (1 -
f)\alpha R_v$，因此 $R_v/R_0 = 1/[f + \alpha(1 - f)]$。取 $f = 0.2$，$\alpha =
1.01$：$1/1.008$，水汽偏轻 \qty{7.9}{\permil}；瑞利模型给出
$0.2^{0.01}$，偏轻 \qty{16.0}{\permil}。每次将凝结物立即移出，
水汽便会偏离得两倍远。
\end{solution}

\begin{solution}{exo:b3:stable-isotopes:11}
$M = m(\delta_s - \Delta)/\Delta = 2000\times(-60 + 2.5)/(-2.5) = \qty{46000}{Gt}$。
于是 $m = 46000\times(-2.5)/(-25 + 2.5) = \qty{5100}{Gt}$，略低于
已发表的 6000 至 \qty{12000}{Gt}，其上限对应较大的
偏移幅度。
\end{solution}

\begin{solution}{exo:b3:stable-isotopes:12}
$\delta_v = (1000 + \delta_{v,0})\exp\int(\alpha - 1)\,\mathrm d\ln f - 1000$，因此
$\mathrm d\delta_v = (1000 + \delta_v)(\alpha - 1)\,\mathrm d\ln f$，且 $f =
e_s/e_{s,0}$ 给出 $\mathrm d\ln f = \mathrm d\ln e_s$。在 \qty{-20}{\celsius} 处，
$\alpha = 1.01872$：$0.01872\times950.9\times0.0960 =
\qty{1.71}{\permil.K^{-1}}$，对应水汽。
\end{solution}

\begin{solution}{pb:b3:stable-isotopes:1}
\textbf{1.} $\alpha^{18}$：1.01605、1.01691、1.01779、1.01872、1.01968、
1.02068；$\alpha^{\mathrm D}$：1.1424、1.1524、1.1630、1.1744、1.1866、1.1997
（温度从 $-5$ 到 \qty{-30}{\celsius}）。

\textbf{2.} $980\times0.425^{0.01605} - 1000 = \qty{-33.4}{\permil}$。

\textbf{3.} $1.01691\times966.6 - 1000 = \qty{-17.0}{\permil}$。

\textbf{4.} 水汽 $966.6\times(0.3976)^{0.01779} - 1000 = \qty{-49.1}{\permil}$；
雪 $1.01872\times950.9 - 1000 = \qty{-31.3}{\permil}$。

\textbf{5.} 水汽 $950.9\times(0.3680)^{0.01968} - 1000 = \qty{-67.6}{\permil}$；
雪 $1.02068\times932.4 - 1000 = \qty{-48.3}{\permil}$。

\textbf{6.} $(932.4/980)\times0.0622 = 0.059$：\qty{5.9}{\percent} 的
氧-18，相对于水的 \qty{6.2}{\percent}：水汽失去重同位素的
速度更快。

\textbf{7.} $(-17.0 + 48.3)/20 = \qty{1.57}{\permil}$ 每度。

\textbf{8.} $t = (-48.3 + 13.6)/0.69 = \qty{-50}{\celsius}$。该关系
适用于地表温度；在冰盖上，地表远比雪形成处的空气寒冷，
而模型忽略了混合与动力学效应。

\textbf{9.} 水汽 $-247.5$、$-352.5$、\qty{-462.7}{\permil}；雪 $\delta\mathrm D
= -132.8$、$-239.6$、\qty{-355.4}{\permil}。

\textbf{10.} $d = -132.8 + 8\times17.02 = 3.3$；随后为 10.7 和
\qty{31.3}{\permil}，过量值向内陆增大。局地斜率
蒸馏的 $\mathrm d\delta\mathrm D/\mathrm d\delta^{18}\mathrm O = (\alpha^{\mathrm D} -
1)(1000 + \delta\mathrm D)/[(\alpha^{18} - 1)(1000 + \delta^{18}\mathrm O)]$，当水汽变得很轻时降至
8 以下（$1000 + \delta\mathrm D$ 的减小速度远快于
$1000 + \delta^{18}\mathrm O$），因此雪的点位于该线之上。

\textbf{11.} 该线给出 $8\times(-48.34) + 10 = \qty{-376.7}{\permil}$，
雪的值比它高 \qty{21}{\permil}。模型只考虑平衡分馏，
而真实的雪在过饱和空气中形成，存在动力学效应，
且气团在途中会与其他气团混合。

\textbf{12.} $\delta_c = (3.25 + 29.99)/0.97001 = 34.27$; $1000\ln\alpha =
33.69$; $T = 18030/66.11 = \qty{272.7}{K}$, \qty{-0.4}{\celsius}.

\textbf{13.} $\delta_c = 36.09$, $1000\ln\alpha = 35.46$, $T = \qty{265.6}{K}$,
\qty{-7.5}{\celsius}：比海水在不凝固时所能达到的温度还要低。
这一变化不可能全部由温度引起。

\textbf{14.} $\delta_w = 125\times0.011 = \qty{+1.375}{\permil}$; $1000\ln(1036.09/1001.375)
= 34.08$；$T = 18030/66.50 = \qty{271.1}{K}$，\qty{-2.0}{\celsius}：降温
\qty{1.6}{\celsius}.

\textbf{15.} $\Delta_{47} = 0.0391\times10^6/272.71^2 + 0.154 = 0.6797$ 与
$0.0391\times10^6/271.12^2 + 0.154 = 0.6859$：相差
\qty{0.006}{\permil}.

\textbf{16.} 它表征重同位素在
矿物中的排布方式，这是一种仅取决于温度的内部平衡。

\textbf{17.} $f_{\mathrm{org}} = (-5 - 0.5)/(-28) = 0.20$.

\textbf{18.} $m = 46000\times(-3.0)/(-60 + 3.0) = \qty{2400}{Gt}$.

\textbf{19.} $46000\times(-3.0)/(-25 + 3.0) = \qty{6300}{Gt}$.

\textbf{20.} $46000\times(-3.0)/(-5 + 3.0) = \qty{69000}{Gt}$，超过
储库本身的量，且发生在地质上极短的时间内，不合理。

\textbf{21.} 岩芯中碳酸盐溶解的深度与范围，它
度量了进入海洋的碳量，且与其组成无关；
结合\omterm{def:b3:stable-isotopes:cie}{碳同位素漂移}的幅度，即可确定 $\delta_s$。

\textbf{22.} 伊尔斯冰穹处的雪的 $\delta^{18}\mathrm O =
\textbf{\qty[text-series-to-math=true]{-48.3}{\permil}}$（\qty{-48.4}{\permil}
，$\alpha$ 连续变化时）。
\end{solution}
